\section{Introduction}\label{sec:introduction}
A Brain-Computer Interface (BCI) is a technology that enables direct communication between the brain and an external device. Electroencephalography (EEG) is one of the most widely used signal sources for BCIs due to its non-invasive nature. In this experiment, EEG signals are acquired with a BrainLink headset, which provides a single dry electrode placed at Fp1 over the prefrontal cortex. The objective is to apply machine learning techniques to these single-channel signals in order to distinguish between three mental states: "Relax", "Focus", and "Stroop".

\noindent The experiment is organized into two parts. In Part 1, each group records its own dataset. In Part 2, these recordings are used to train and evaluate classifiers.

\noindent Our task is to design an effective preprocessing and feature extraction pipeline, to implement and compare at least four classifiers (MLP, Support Vector Machine, Random Forest, and one self-selected model), and to apply post-processing strategies, with the goal of raising the overall mean accuracy of the best model to 75\% or higher. All models are assessed under Leave-One-Subject-Out (LOSO) cross-validation.

\noindent The three states are expected to leave different traces at Fp1. Relaxed wakefulness is associated with stronger alpha activity, while sustained attention and cognitive conflict, such as the Stroop task, are generally expected to show a shift toward faster activity and more frontal theta. The Stroop task in this experiment is answered with key presses, so movement and muscle artifacts also add amplitude cues. In the following sections we describe how our pipeline exploits these cues and how well it generalizes across the three group members.
